% Insertable contribution: no documentclass, preamble, or bibliography command.
% Requires amsmath, amssymb, booktabs, graphicx, and natbib.
% Set \goshaRoot to this directory relative to the main compilation directory.
\providecommand{\goshaRoot}{report/gosha}
\input{\goshaRoot/numbers.tex}
\input{\goshaRoot/tables/main_results.tex}
\input{\goshaRoot/tables/tie_seed_summary.tex}

\section{Evaluation Protocol}
\label{sec:gosha-protocol}
We audit two frozen FB15k-237 checkpoints without test-based selection.
Each test fact produces head $(?,r,t)$ and tail $(h,r,?)$ queries against
all 14,541 entity IDs: $N=40{,}932$ queries, 20,466 per mode. Canonical order
is test-file order, all head queries then all tail queries. We exclude
other known true answers from
$\mathcal{D}_{\mathrm{train}}\cup\mathcal{D}_{\mathrm{valid}}
\cup\mathcal{D}_{\mathrm{test}}$, while always retaining the target.
Using test facts for benchmark filtering does not permit training on them.

Let $K_q$ be retained candidates and $y_q$ the target. Higher raw scores are
better. We use finite, unrounded float32 scores without AMP/TF32 and read
$s_q(y_q)$ from the \emph{same candidate row}, not a separate forward pass:
\begin{align}
 B_q&=\sum_{e\in K_q}\mathbf{1}[s_q(e)>s_q(y_q)],\\
 T_q&=\sum_{e\in K_q}\mathbf{1}[s_q(e)=s_q(y_q)].
\end{align}
Equality is exact, without tolerance or rounding; $T_q$ includes the target.
A target tie has $T_q>1$. Optimistic and pessimistic ranks are
$r_q^{\mathrm{O}}=B_q+1$ and $r_q^{\mathrm{P}}=B_q+T_q$.
For a uniform position inside the target's tie group,
\begin{align}
 \mathbb{E}[\mathrm{RR}_q]&=\frac{1}{T_q}
       \sum_{j=1}^{T_q}\frac{1}{B_q+j},\label{eq:gosha-rr}\\
 \mathbb{E}[\mathrm{H@}k_q]&=
       \frac{\min(\max(k-B_q,0),T_q)}{T_q}.
\end{align}
The expectation in Eq.~\ref{eq:gosha-rr} averages reciprocals; it is not
$1/\mathbb{E}[r_q]$. For a deterministic rank policy,
$\mathrm{MRR}=N^{-1}\sum_q1/r_q$ and
$\mathrm{H@}k=N^{-1}\sum_q\mathbf{1}[r_q\leq k]$, for $k\in\{1,3,10\}$.
Expected-random metrics average the corresponding expectations.
Our primary effect is
$\Delta=\mathrm{MRR}_{\mathrm{O}}-\mathbb{E}[\mathrm{MRR}_{\mathrm{random}}]$.

Each candidate row is computed once; saved $B_q,T_q$ serve all policies.
For each evaluation seed $a\in\{0,1,2,3,4\}$, we separately draw
$r_q^{(a)}=B_q+U_q^{(a)}$, with $U_q^{(a)}$ uniform on
$\{1,\ldots,T_q\}$, using a fresh NumPy \texttt{default\_rng(a)} in canonical
query order. We report all five values, their mean and sample standard
deviation ($\mathrm{ddof}=1$). These draws vary evaluation at fixed weights,
not training; they do not estimate training uncertainty.

\section{Results}
\label{sec:gosha-results}
Table~\ref{tab:gosha-main} gives expected-random MRR of
\GoshaRotateMRR{} for RotatE and \GoshaConvKBMRR{} for corrected ConvKB.
The nonzero gaps are \GoshaRotateGap{} and \GoshaConvKBGap{} MRR units.
All three policies give identical Hits@1/3/10. Only
\GoshaRotateTies{} and \GoshaConvKBTies{} queries have target ties;
every non-singleton target group has $T_q=2$.
Figure~\ref{fig:gosha-ties} displays these rare events directly, rather than
hiding them beside a dominant no-tie bar. Table~\ref{tab:gosha-ties}
summarizes the five draws; their means need not equal analytic expectation
exactly. Head/tail results, exact histograms and all seeds are provided
separately as machine-readable tables.

\begin{figure}[t]
 \centering
 \ifdim\linewidth>3.2in
  \includegraphics[width=3.2in]{\goshaRoot/figures/target_ties.pdf}
 \else
  \includegraphics[width=\linewidth]{\goshaRoot/figures/target_ties.pdf}
 \fi
 \caption{Queries with exact target ties after filtering. Each denominator is
 40,932 queries; the vertical axis is a linear percentage scale.}
 \label{fig:gosha-ties}
\end{figure}

\section{Reproducibility}
\label{sec:gosha-repro}
The full scoring run used a Tesla T4, Python 3.13.16, PyTorch
2.11.0+cu130 and NumPy 2.3.3. Query batch/candidate chunk sizes were
8/2048 for RotatE and 2/512 for ConvKB. Metadata records numerical backend
settings, checkpoint/archive/data/source hashes and the ordered query IDs.
Code and results are versioned in the repository
\citep{audit-repository}.

We distinguish \emph{reconstructing saved results} from \emph{scoring models
again}. The former passed in a fresh clone and an isolated environment
without PyTorch/GPU. Independent direct reciprocal sums check every policy
and draw against JSON/CSV/Markdown. Benchmark-aware checks verify target IDs
and retained counts against restored splits. In a second isolated environment,
tests and validation-only CPU scoring of both checkpoints passed. A new full
GPU rerun remains blocked locally by unavailable CUDA. Frozen packages are
distributed separately; no public weight download has been verified.
Commands, locks, exit codes and logs accompany this contribution.

\section{Interpretation and Limitations}
\label{sec:gosha-discussion}
\citet{sun-etal-2020-evaluation}, Table~2, report ConvKB
TOP/RANDOM/BOTTOM MRR of 0.407/0.243/0.130, giving a displayed TOP--RANDOM
gap of 0.164. Their RANDOM is a mean over five runs, whereas our primary
RANDOM is analytic; we retain the five draws separately. Their RotatE values
all round to 0.336; its unrounded gap remains unresolved.

On our selected budget checkpoints, tie-breaking has negligible impact;
the paper's large ConvKB effect is not reproduced. This does not refute the
paper: training recipes, implementation, validation selection and numerical
environments differ, without causal ablations. ConvKB has corrected h/r/t
mixing but does not reproduce every upstream training choice. One checkpoint
per model cannot quantify training variability. Float32 ties may change
across devices or batching. Precision, weights, seeds and hyperparameters
were not tuned to these results.
